\section{Related Work}

\subsection{Slovene language resources and social-media corpora}

Written standard Slovene is well served by reference resources. Gigafida 2.0 \parencite{krek2020b-gf,gigafida-clarin} is the main reference corpus, with about 1.1 billion words of written standard Slovene; ssj500k \parencite{krek2020a-ssj500k} provides manually annotated training data; siParl 3.0 \parencite{parliamentaryCorpus} covers parliamentary proceedings from 1990 to 2022. The broader methodological context is surveyed in \textcite{studije2005korpusnem}.

The most directly relevant prior work is the JANES project (\textit{Jezikovna analiza nestandardne slovenščine}), a central Slovene resource for user-generated content across social media and the web \parencite{erjavec2019janes}. JANES mapped out much of the territory: the feasibility of building Slovene social-media corpora and the specific challenges they pose, including non-standard spelling and code-switching. The present work is much smaller in scale and ambition; it revisits some of the same questions in the context of Bluesky.

The closest language-specific Bluesky precedent located so far is \textcite{tjongKimSang2025bluesky}, who evaluates Bluesky as a successor data source for Dutch social-media research after Twitter/X API restrictions made the TwiNL Twitter archive difficult to extend. That study collected 4.2 million Dutch Bluesky posts from November 2024 to May 2025 through keyword-based search and demonstrates compact corpus-linguistic analyses such as word-frequency comparison, daily rhythms, event response, hashtags, and political vocabulary. The present paper is narrower in language-community size, but differs methodologically by using protocol-level ATProto backfill from discovered authors and by making post-level language validation a central part of corpus construction.

\subsection{Bluesky and ATProto}

Prior work characterises Bluesky as both Twitter-like and decentralised, and links parts of its early growth to events around Twitter/X \parencite{failla2024bluesky,quelle2025bluesky}. Platform access conditions also matter here: Bluesky operated as an invite-only beta through 2023, regulating growth through waitlists and invite codes \parencite{bsky-beta-2023,bsky-beta-roadmap-2023}, and only opened to the public on 6~February 2024, when invites were removed \parencite{bsky-public-2024}. The present corpus cannot establish why individual Slovene users joined Bluesky, but the late-2024 surge visible in Section~\ref{sec:empirical} should be interpreted against this broader platform context: public opening had already removed one barrier to entry, while the Slovene corpus shows its main surge only later, suggesting that public opening was a background enabling condition rather than a sufficient explanation on its own.

Technically, Bluesky is built on the AT Protocol (ATProto), a federated specification that separates identity from hosting: user identities are DIDs (decentralised identifiers), and public records are stored in per-account repositories on Personal Data Servers \parencite{kleppmann2024atproto,atproto-spec,atproto-repository}. ATProto also provides public synchronisation mechanisms: repository exports and a repository event stream, commonly called the firehose \parencite{atproto-sync,atproto-streaming}. From a corpus-collection standpoint, this means that public posts can be collected through protocol-level mechanisms rather than through a separate academic access programme. DID resolution allows a collector to follow an account regardless of which PDS hosts it, a property exploited by the multi-PDS collection approach described in this paper.

Existing empirical work on Bluesky has mostly operated at large scale and focused on platform-wide, network-level, or topic-specific phenomena. \textcite{failla2024bluesky} assembled a dataset of about 235 million posts from over four million users for behavioural analysis and reported large language-tag clusters, especially English, Japanese, and German. \textcite{quelle2025bluesky} used Bluesky data to characterise network topology, political polarisation, and user-generated feeds. More recent work treats Bluesky's discovery mechanisms as research objects in their own right: \textcite{balduf2025starterpacks} analyse starter packs as a mechanism for social bootstrapping, while \textcite{hanny2026disaster} use multilingual keyword-based Bluesky collection for disaster analysis. These studies make clear that Bluesky is not simply a Twitter replacement: feeds, starter packs, repository access, and language metadata all affect what can be collected and interpreted.

Language-specific Bluesky corpus work appears sparse in the literature reviewed here. The Dutch study by \textcite{tjongKimSang2025bluesky} is therefore a useful comparison point, while this paper instead offers a smaller but more validation-oriented case study for Slovene and for ATProto-based collection of a small language community.

\subsection{Language identification in corpus construction}

Building a language-specific corpus from a multilingual stream typically combines two signals: platform-supplied language metadata and automatic language identification. Neither is reliable on its own. Bluesky posts may contain a \texttt{langs} field with one or more BCP-47 language codes, and the official client auto-detects post languages while allowing users to override them \parencite{bsky-posts}. Large-scale Bluesky studies have also found that language metadata can be inconsistent or incomplete: \textcite{failla2024bluesky} report non-standard and ill-formed language tags, while \textcite{hanny2026disaster} caution that language metadata was inaccurate enough in a multilingual disaster crawl that filtering only by existing tags would risk missing relevant posts.

Automatic language identification tools---langid.py \parencite{lui2012langid} and langdetect \parencite{nakatani2010langdetect} are used here---are also imperfect, especially for short texts and closely related South Slavic languages. Using consensus across multiple detectors and requiring agreement with author-supplied tags improves precision, which is the rationale for the tiered inclusion strategy described in Section~\ref{sec:corpus-construction}.
